\section{Exposure in experiments that switch at regenerations}\label{sec:exposure}

A switchback experiment assigns one of two treatments, the arms, to a whole system for a
period of time, then the other, and alternates. The quantity under the experimenter's
control is the exposure, the number of time slots spent in the treatment arm. The designs
in the literature switch at fixed times or at the hits of a recurrent
state~\cite{bsz2023,huwager2022,gjr2020,ljkw2023,jky2023,hu2026}, and each of them assumes
that the carryover from one period to the next decays geometrically, through a Markov
state on a finite set, a bounded queue, or a fixed memory. Then the regeneration cycles
have exponential moments and the exposure count is Gaussian at scale $\sqrt N$. When the
system is a queue with heavy-tailed service the cycles have infinite variance, and none of
these results applies. This section works out the exposure law for a switchback that
switches arm at the regenerations of a slotted queue, where the count is exactly the
occupation count $A_N$ of this paper and the cycles are the runs. The results are the
theorems of Sections~\ref{sec:bulk} to~\ref{sec:shape} read in the language of the design,
once the cycle law is shown to satisfy the hypotheses. The second step is the content of
Proposition~\ref{prop:cycle}. Its first part, the local tail of the cycle, is known. The
difference bounds (LD) and (LD2) for the cycle are new, and they are what the mode count
needs.

\subsection{The queue, the design and the reduction}\label{sec:exposure-model}

Time is slotted. In each slot a customer arrives with probability $\lambda\in(0,1)$, a
constant unrelated to the factor $\lambda_N$ of~\eqref{eq:lambdaN}, and brings $Y$ slots of
work, where the $Y$ are independent with values in $\{1,2,\dots\}$. So the work $W$ that
arrives in a slot has $\Prob(W=0)=1-\lambda$ and $\Prob(W=n)=\lambda f_Y(n)$ for $n\ge1$, and
its mean is $\rho=\lambda\E Y$, which we assume to be less than $1$. The server does one unit
of work per slot, so the workload at the end of slot $t$ is
$V_t=\max(V_{t-1}-1,0)+W_t$. A regeneration is a slot that ends with $V_t=0$, and the cycle
$C$ is the number of slots between consecutive regenerations. An idle slot is a cycle of
length $1$, so $\Prob(C=1)=1-\lambda$. Since the fraction of busy slots is $\rho$ and the
regenerations are the idle slots, $\E C=1/(1-\rho)$.

The 2 arms set their own arrival probabilities $\lambda_\plus$ and $\lambda_\minus$ and
share the law of $Y$, so their loads $\rho_\pm=\lambda_\pm\E Y$ differ. The experiment
switches arm at every regeneration. At a switch the system is empty, so the cycles are
independent, each cycle has the cycle law of its own arm, and consecutive cycles belong to
opposite arms. Hence the number of slots spent in the treatment arm among the first $N$ is
the count $A_N$ of Section~\ref{sec:runs}, with the plus runs the cycles under treatment and
the minus runs the cycles under control. In this section $f_\pm$ and $T_\pm$ are the cycle
laws of the two arms, $\mu_\pm=1/(1-\rho_\pm)$, $s=\mu_\plus+\mu_\minus$ and $p=\mu_\plus/s$ is the
long-run treated share. Three starts occur. A fresh start of either arm is an experiment
that starts on an empty system. A random-time start is an experiment that starts at a
typical time of a long-running switchback, and its initial run is the stationary initial run
$\nu_\sgn=T_\sgn/s$ of Section~\ref{sec:prep-stationary}. An aged start is an experiment that
starts $t_a$ slots after a switchback started on an empty system, and its law is the aged law
of Section~\ref{sec:aging}.

\subsection{The cycle law}\label{sec:exposure-cycle}

Let $U_n=W_1+\dots+W_n$. Starting from an empty system, $V_1=W_1$ and $V_t=V_{t-1}-1+W_t$ as
long as $V_{t-1}\ge1$, so $V_t=U_t-t+1$ up to the end of the cycle, and
$C=\inf\{t\ge1:U_t-t=-1\}$ is the first passage of the downward skip-free walk $U_t-t$ to
$-1$. The hitting time theorem of Kemperman~\cite{kemperman1961} (see van der Hofstad and
Keane~\cite{vdhk2008} for a short proof) gives the cyclic lemma
\begin{equation}\label{eq:cyclic}
 \Prob(C=n)=\frac1n\,\Prob(U_n=n-1)\qquad(n\ge1).
\end{equation}
Equivalently, $C$ is the total progeny of a Galton--Watson tree with offspring law
$W$~\cite{dwass1969}. Let $W'=W+1$, a law on $\{1,2,\dots\}$ with mass function
$f_{W'}(n)=f_W(n-1)$ and mean $\mu'=1+\rho$, and let $V_m=f_{W'}^{*m}$. Then $U_n=U'_n-n$ with
$U'_n$ a sum of $n$ copies of $W'$, so~\eqref{eq:cyclic} reads
\begin{equation}\label{eq:phi}
 \varphi(n):=n\,f_C(n)=V_n(2n-1),
\end{equation}
a convolution power evaluated at the deviation $L=(2n-1)-n\mu'=(1-\rho)n-1$ above its mean.
This puts the cycle law inside the big-jump range of the lemmas of Section~\ref{sec:shape}
and Appendix~\ref{app:modes}, with the number of summands $m=n$ of the order of the deviation.

\begin{proposition}[the cycle law]\label{prop:cycle}
Let $Y$ satisfy (LR) with index $\alpha\in(1,2)$, that is $f_Y(n)\sim\alpha T_Y(n)/n$ with
$T_Y\in\RV_{-\alpha}$.
\begin{enumerate}[(a)]
\item $f_C(n)\sim f_W\bigl((1-\rho)n\bigr)\sim\lambda(1-\rho)^{-\alpha-1}f_Y(n)$ and
$T_C(n)\sim\lambda(1-\rho)^{-\alpha-1}T_Y(n)$. So $C$ satisfies (LR) with the index $\alpha$,
and (SPAN).
\item If $Y$ also satisfies (LD2), then $C$ satisfies (LD) and (LD2), and
$f_C(n+1)-f_C(n)=-(\alpha+1+o(1))f_C(n)/n$.
\item If 2 arms share $Y$ and have loads $\rho_\pm<1$, their cycle laws satisfy (EQ), and
$T_\plus(n)/T_\minus(n)\to(\lambda_\plus/\lambda_\minus)\bigl((1-\rho_\minus)/(1-\rho_\plus)\bigr)^{\alpha+1}$.
\end{enumerate}
\end{proposition}

Part (a) is known. For the total progeny of a Galton--Watson tree it is
Kortchemski~\cite[Eq.~(14)]{kortchemski2015} and Janson~\cite[Remark~18.34]{janson2012}, and
for the busy period of the M/G/1 queue in continuous time it is Denisov and
Shneer~\cite[Corollary~2.2]{denisovshneer2010}, whose $f_\tau(t)\sim f_B((1-\rho)t)$ is the
analogue of the first relation in (a). The global tail of the busy period, with the factor
$\E\nu=1/(1-\rho)$ for the mean number served, is due to De Meyer and
Teugels~\cite{demeyerteugels1980} and Zwart~\cite{zwart2001}. We include the short proof of (a)
because (b) continues it. We did not find (b) in the literature.

\begin{proof}
(a) The law $W'$ satisfies (LR) with index $\alpha$, since $f_{W'}(n)=\lambda f_Y(n-1)$ for
$n\ge2$, and (SPAN), since $f_{W'}(1)=1-\lambda>0$ and $f_{W'}(n)>0$ for all large $n$. The
local big-jump theorem of Denisov, Dieker and Shneer~\cite[Theorem~9.1]{dds2008}, whose
hypotheses were verified for a law with (LR) in the proof of Lemma~\ref{lem:local-bigjump},
gives for every $\varepsilon>0$ an $r_0$ with
$|\Prob(U'_r-r\mu'\in(x,x+1])/(r\Prob(x<W'-\mu'\le x+1))-1|\le\varepsilon$ for all $r\ge r_0$
and $x\ge x_r=(1-\rho)r/2$, because $rT_{W'}(x_r)\to0$. Take $r=n$ and
$x=2n-1-n\mu'-\frac12$. The event is $\{U'_n=2n-1\}$ and
$\Prob(x<W'-\mu'\le x+1)=f_{W'}(n')$ with $n'=(1-\rho)n+O(1)$. By~\eqref{eq:f-uct} for the
law $W'$, $f_{W'}(n')=f_{W'}((1-\rho)n)(1+o(1))$, and $f_{W'}((1-\rho)n)=f_W((1-\rho)n-1)\sim f_W((1-\rho)n)$
by the same lemma. Hence $\varphi(n)=nf_W((1-\rho)n)(1+o(1))$, which is the first relation,
and the second is the regular variation of $f_W=\lambda f_Y$ on $n\ge1$. The tail follows by
Karamata's theorem, and (LR) for $C$ follows from (LR) for $Y$. Finally $f_C(1)=1-\lambda>0$ and
$f_C(n)>0$ for large $n$, which gives (SPAN).

(b) The law $W'$ satisfies (LD2) because $Y$ does, and (LD) by Lemma~\ref{lem:ld2}. Write
$f=f_{W'}$, $a=a_{W'}$, and $L_x=x-n\mu'$ for the deviation of the point $x$ at $m=n$
summands, so that $L_{2n-1}=L$ and $L_{2n+i}=L+i+1$. For large $n$, $L\ge Ka(n)$ for any fixed
$K$, since $a(n)=o(n)$. By~\eqref{eq:phi},
\[
 \varphi(n+1)-\varphi(n)=V_{n+1}(2n+1)-V_n(2n-1)
 =\delta V_n(2n+1)-\nabla V_n(2n+1)-\nabla V_n(2n),
\]
in the notation of Lemma~\ref{lem:diff-bigjump}. That lemma, applied to the law $W'$, gives,
with $o(1)$ as $n\to\infty$,
\begin{align*}
 \delta V_n(2n+1)&=f(L)+\mu'(\alpha+1)\frac{nf(L)}L+o\Bigl(f(L)+\frac{nf(L)}L\Bigr),\\
 \nabla V_n(2n+i)&=(\alpha+1)\frac{nf(L)}L\bigl(1+o(1)\bigr)\qquad(i=0,1),
\end{align*}
where we used $f(L+i)=f(L)(1+O(1/L))$ for $|i|\le2$, by (LD). Since $n/L\to1/(1-\rho)$ and
$\mu'-2=-(1-\rho)$,
\[
 \varphi(n+1)-\varphi(n)=f(L)\Bigl[1+(\alpha+1)\frac nL(\mu'-2)+o(1)\Bigr]=-\alpha f(L)(1+o(1)).
\]
By (a), $\varphi(n)=nf(L)(1+o(1))$. Hence
\[
 f_C(n+1)-f_C(n)=\frac{\varphi(n+1)-\varphi(n)}{n+1}-\frac{\varphi(n)}{n(n+1)}
 =-\frac{(\alpha+1+o(1))f(L)}{n},
\]
and since $f_C(n)\sim f(L)$ by (a), this is (LD) for $C$ with the stated constant. For (LD2) we use the second-difference bounds
of Lemma~\ref{lem:cycle-second} below, with $c=(1-\rho)/2$, which apply at $m=n$ and at the
points $2n-1\le x\le2n+3$ for large $n$. Put $D(n)=\varphi(n+1)-\varphi(n)$. Then
\begin{align*}
 D(n+1)-D(n)={}&\bigl[\delta V_{n+1}(2n+3)-\delta V_n(2n+1)\bigr]
 -\bigl[\nabla V_{n+1}(2n+3)-\nabla V_n(2n+1)\bigr]\\
 &-\bigl[\nabla V_{n+1}(2n+2)-\nabla V_n(2n)\bigr]\\
 ={}&\delta^2V_n(2n+3)-\nabla\delta V_n(2n+3)-\nabla\delta V_n(2n+2)\\
 &-\nabla\delta V_n(2n+3)+\nabla^2V_n(2n+3)+\nabla^2V_n(2n+2)\\
 &-\nabla\delta V_n(2n+2)+\nabla^2V_n(2n+2)+\nabla^2V_n(2n+1),
\end{align*}
where each bracket was split into a difference in $m$ at the larger point and a difference in
$x$ at $m=n$, and $X(x+2)-X(x)=-\nabla X(x+2)-\nabla X(x+1)$. By Lemma~\ref{lem:cycle-second}
each of the nine terms is at most $Cf(L)/L$ in absolute value, since $n\le L/c$. So
$|\varphi(n+2)-2\varphi(n+1)+\varphi(n)|\le Cf(L)/n$. With $g(n)=1/n$ the product rule gives
\[
 \Delta^2(g\varphi)(n)=g(n+2)\Delta^2\varphi(n)+\Delta\varphi(n)\bigl[\Delta g(n)+\Delta g(n+1)\bigr]+\varphi(n+1)\Delta^2g(n),
\]
and the three terms are $O(f(L)/n^2)$, $O(f(L))\cdot O(n^{-2})$ and $O(nf(L))\cdot O(n^{-3})$. Since
$f(L)\asymp f_C(n)$ by (a), this is (LD2) for $C$.

(c) By (a), $T_\pm(n)\sim\lambda_\pm(1-\rho_\pm)^{-\alpha-1}T_Y(n)$, and the ratio of the two
constants is the stated limit.
\end{proof}

The lemma used for (LD2) is the second-difference bound of Lemma~\ref{lem:second-bigjump},
extended from deviations beyond the mean, $L\ge m\mu$, to deviations of the order of the
number of summands, together with its consequences for the differences in $m$. The proof of
Lemma~\ref{lem:second-bigjump} uses $L\ge m\mu$ at two places only, to bound
$mT(\delta L)$ by $C_\delta LT(L)$ and $m$ by $CL$, and both bounds hold under $L\ge cm$ with
constants that depend on $c$.

\begin{lemma}[second differences at deviations of order $m$]\label{lem:cycle-second}
Let $f$ be a law on $\{1,2,\dots\}$ with mean $\mu$ that satisfies (SPAN), (LR) and (LD2), let
$a$ be its scale function and $V_m=f^{*m}$, and fix $c>0$. There are $K_2$, $L_1$ and $C$,
depending on $f$ and $c$, such that for all $m\ge1$ and all integers $n$ with
$L=n-m\mu\ge\max(K_2a(m),L_1,cm)$,
\[
 |\nabla^2V_m(n)|\le C\,m\,\frac{f(L)}{L^2},\qquad
 |\nabla\delta V_m(n)|\le C\,\frac{f(L)}L,\qquad
 |\delta^2V_m(n)|\le C\,\frac{f(L)}L .
\]
\end{lemma}

\begin{proof}
We write $T$ for the tail of $f$. Throughout, $L\ge cm$ gives $m\le L/c$, $mf(L)\le CT(L)/c$
by (LR), and $mT(\delta L)\le C_{\delta,c}LT(L)$ by Lemma~\ref{lem:rv}(ii), and
Lemma~\ref{lem:rv}(ii) also gives $f(L')\le C_bf(L)$ for $L'\ge bL$ and $L$ large. By
Lemma~\ref{lem:diff-bigjump} there is $K$ such that $|\nabla V_m(x)|\le Cmf(L_x)/L_x$,
$V_m(x)\le Cmf(L_x)$ and $|\delta V_m(x)|\le C(f(L_x)+mf(L_x)/L_x)$ whenever $L_x=x-m\mu\ge Ka(m)$.

\emph{The first bound.} This is Lemma~\ref{lem:second-bigjump} under $L\ge cm$ in place of
$L\ge m\mu$. In its proof the hypothesis $L\ge m\mu$ is used to get $mT(\delta L)\le C_\delta LT(L)$,
which makes $mT(w)\le1$ for $w=\lceil\delta L\rceil$ and $L\ge L_0$, and to get
$mLf(L)\le CmT(L)\le C$ and $mT(L)L/a(m)\le C$. The first bound holds under $L\ge cm$ with
$C_{\delta,c}$, and the other two use only $L\ge K_1a(m)$ through~\eqref{eq:potter-m}. Every
other step of that proof, and of the proofs of Lemmas~\ref{lem:trunc-local},
\ref{lem:truncated} and~\ref{lem:diff-bigjump} that it invokes, is stated for $L\ge Ka(m)$
with no upper bound on $m$. So the first bound holds for $L\ge\max(K_1a(m),L_0,cm)$, with
constants depending on $c$.

\emph{The second bound.} Since $\delta V_m=f*V_m-V_m$ and $\nabla$ commutes with convolution,
\[
 \nabla\delta V_m(n)=\sum_{x\le L/2}f(x)\bigl[\nabla V_m(n-x)-\nabla V_m(n)\bigr]
 -\nabla V_m(n)\,T(L/2)+\sum_{y<n-L/2}f(n-y)\,\nabla V_m(y).
\]
For $x\le L/2$ and $1\le i\le x$ the point $n-i+1$ has deviation at least $L/2$, so the first
bound at the constant $c/2$ gives
$|\nabla V_m(n-x)-\nabla V_m(n)|\le\sum_{i=1}^{x}|\nabla^2V_m(n-i+1)|\le Cxmf(L)/L^2$, and the first
sum is at most $C\mu mf(L)/L^2\le Cf(L)/L$. The middle term is at most
$Cmf(L)T(L/2)/L\le C(mT(L))f(L)/L\le Cf(L)/L$. In the last sum put $Y=\lceil n-L/2\rceil-1$ and sum
by parts,
\[
 \sum_{y\le Y}f(n-y)\bigl[V_m(y-1)-V_m(y)\bigr]
 =\sum_{y\le Y-1}\bigl[f(n-y-1)-f(n-y)\bigr]V_m(y)-f(n-Y)V_m(Y).
\]
For $y\le Y$ we have $n-y\ge L/2$, so $f(n-y)\le Cf(L)$ and, by (LD),
$|f(n-y-1)-f(n-y)|\le Cf(L)/L$. Since $\sum_yV_m(y)\le1$, the sum is at most $Cf(L)/L$. The
boundary term has $Y-m\mu\ge L/2-1$, so $V_m(Y)\le Cmf(L)$ and the term is at most
$Cmf(L)^2\le C(mT(L)/L)f(L)\le Cf(L)/L$.

\emph{The third bound.} Since $\delta^2V_m=f*\delta V_m-\delta V_m$,
\[
 \delta^2V_m(n)=\sum_{x\le L/2}f(x)\bigl[\delta V_m(n-x)-\delta V_m(n)\bigr]
 -\delta V_m(n)\,T(L/2)+\sum_{y<n-L/2}f(n-y)\,\delta V_m(y).
\]
By the second bound at the constant $c/2$, the first sum is at most
$\sum_xf(x)\,x\,Cf(L)/L\le C\mu f(L)/L$. The middle term is at most
$C(f(L)+mf(L)/L)T(L)\le Cf(L)T(L)=o(f(L)/L)$, since $LT(L)\to0$. For the last sum put
$\gamma(y)=f(n-y)\ones_{\{y<n-L/2\}}$. With $S_m$ of law $V_m$ and $X$ an independent copy of the
law $f$, $\delta V_m(y)=\Prob(S_m+X=y)-\Prob(S_m=y)$, so the last sum is
$\E[\gamma(S_m+X)-\gamma(S_m)]$. If $S_m+X<n-L/2$, then both arguments lie below $n-L/2$, and
(LD) applied $X$ times gives $|\gamma(S_m+X)-\gamma(S_m)|\le CXf(L)/L$, with expectation at most
$C\mu f(L)/L$. If $S_m<n-L/2\le S_m+X$, the difference is $\gamma(S_m)\le Cf(L)$, and this event
has probability $\sum_{k\ge1}T(k)V_m(n'-k)$ with $n'=\lceil n-L/2\rceil$. For $k\le L/4$ the
point $n'-k$ has deviation at least $L/4-1$, so $V_m(n'-k)\le Cmf(L)$, and these terms add up
to at most $C\mu mf(L)\le CT(L)$. For $k>L/4$, $T(k)\le T(L/4)\le CT(L)$, and the $V_m$ add up
to at most $1$. So the event has probability at most $CT(L)$, and its contribution is at most
$Cf(L)T(L)=o(f(L)/L)$. If $S_m\ge n-L/2$, both values of $\gamma$ vanish. Enlarging $K_2$ and
$L_1$ so that every deviation used above lies in the range of the earlier bounds completes
the proof.
\end{proof}

\subsection{The exposure law}\label{sec:exposure-law}

Throughout this subsection $Y$ satisfies (LR) with index $\alpha\in(1,2)$, so both cycle laws
satisfy (RV), (LR), (SPAN) and (EQ) by Proposition~\ref{prop:cycle}, and the exposure count
$A_N$ is the count of this paper with these two run laws. By Proposition~\ref{prop:cycle}(c)
and Lemma~\ref{lem:rv}, the ratio $\theta_{\sgn,N}=b_{\sgn,N}/B_N$ of the two scales converges,
to $\theta_\sgn=\min\{1,(\kappa_\sgn/\kappa_{\osgn})^{1/\alpha}\}$ with
$\kappa_\sgn=\lambda_\sgn(1-\rho_\sgn)^{-\alpha-1}$, so the moving profile $h_N$ converges
uniformly to the fixed density $h_\theta$ of~\eqref{eq:profile}.

\begin{corollary}[the exposure imbalance]\label{cor:exposure-bulk}
For a fresh start of either arm,
\[
 \sup_{0\le k\le N}\Bigl|B_N\Prob(A_N=k)-h_\theta\Bigl(\frac{k-pN}{B_N}\Bigr)\Bigr|\to0,\qquad
 \sup_{N,k}B_N\Prob(A_N=k)<\infty .
\]
For a random-time start the same holds over $1\le k<N$. It holds over $0\le k\le N$ when
$\alpha>\varphi$, and when $\alpha<\varphi$ it fails at the endpoints, where
$B_N\Prob(A_N\in\{0,N\})\to\infty$. In all cases $(A_N-pN)/B_N$ converges in law to
$W_\theta$, so the treated share $A_N/N$ misses $p$ by a random error of exact order
$N^{1/\alpha-1}$ with a skewed stable law, in addition to the deterministic imbalance
$(2p-1)N$ of the counts when $\rho_\plus\ne\rho_\minus$.
\end{corollary}

\begin{proof}
The fresh start is Theorem~\ref{thm:bulk}(a) and~(b) with $h_N\to h_\theta$. The random-time
start satisfies the bulk criterion by Corollary~\ref{cor:stationary}, so
Theorem~\ref{thm:preparation}(b) gives the statement over $1\le k<N$. Over the whole range it
holds if and only if $B_NQ_\sgn(N)\to0$ for both arms, and
$B_NQ_\sgn(N)\in\RV_{1/\alpha+1-\alpha}$ by Corollary~\ref{cor:stationary}(a), which tends to
zero for $\alpha>\varphi$ and to infinity for $\alpha<\varphi$. The convergence in law follows
from the local statement over $1\le k<N$ and $Q_\sgn(N)\to0$.
\end{proof}

The limit in law is classical. For the sojourn time of an alternating renewal process it is
due to Tak\'acs~\cite{takacs1957}, and Godr\`eche and Luck~\cite{godrecheluck2001} give it for
one law. What the corollary adds is the local form at every $k$, uniformly, and the bound
$C/B_N$ at every point, together with the exceptions at the endpoints.

\begin{corollary}[severe imbalance and no switch]\label{cor:exposure-tail}
Fix $\varepsilon\in(0,1-p)$ and put $\ell_\varepsilon=1-\varepsilon/(1-p)$.
\begin{enumerate}[(a)]
\item For a fresh start of either arm,
\[
 \Prob(N-A_N<\varepsilon N)\sim NT_\plus(N)\,\frac{\varepsilon}{\mu_\minus}\,\ell_\varepsilon^{-\alpha},
\]
and for a random-time start
\[
 \Prob(N-A_N<\varepsilon N)\sim NT_\plus(N)\Bigl[\frac{\varepsilon}{\mu_\minus}\,\ell_\varepsilon^{-\alpha}
 +\frac{\ell_\varepsilon^{1-\alpha}}{s(\alpha-1)}\Bigr].
\]
The mirror statements hold for $\Prob(A_N<\varepsilon N)$ with the signs exchanged.
\item (No switch.) The event $\{A_N\in\{0,N\}\}$ is the event that the first cycle covers all
$N$ slots. For a fresh start in the treatment arm its probability is $T_\plus(N)$ exactly. For
a random-time start it is $Q_\plus(N)+Q_\minus(N)$ exactly, with
$Q_\sgn(N)=s^{-1}\sum_{d\ge N}T_\sgn(d)\sim NT_\sgn(N)/(s(\alpha-1))$. Hence
\[
 \frac{\Prob_{\mathrm{random}}(\text{no switch})}{\Prob_{\mathrm{fresh}\,\plus}(\text{no switch})}
 \sim\frac{N}{s(\alpha-1)}\Bigl(1+\frac{T_\minus(N)}{T_\plus(N)}\Bigr).
\]
For an aged start with $t_a/N\to\theta\in(0,\infty)$ the ratio carries the factor
$1-(1+\theta)^{1-\alpha}$ in the first term, by Theorem~\ref{thm:aging}(b).
\end{enumerate}
\end{corollary}

\begin{proof}
(b) is Proposition~\ref{prop:fresh-identity}, Proposition~\ref{prop:first-run} and
Corollary~\ref{cor:stationary}(a). For (a), let $\delta\in(0,\varepsilon)$. The points
$(1-\varepsilon)N\le k\le(1-\delta)N$ have $L_k\ge(\ell_\varepsilon N)(1+o(1))$, so they lie in
the zone of Theorem~\ref{thm:bigjump}(b) for every $K$ once $N$ is large, because
$\lambda_NB_N=o(N)$, and there $\widehat G_N(k)\le C_\varepsilon T_\plus(N)$. By the additive form
of Theorem~\ref{thm:bigjump}(b), the sum of $P_N(k)$ over these points is
$\sum_kG_N(k)+o(NT_\plus(N))$. By the computation in the proof of
Corollary~\ref{cor:infinite-density}, together with the uniform convergence of
Lemma~\ref{lem:rv}(i) and~\eqref{eq:gamma-asymp}, $G_N(\lfloor vN\rfloor)/T_\plus(N)\to I_\plus(v)$
uniformly in $v\in[1-\varepsilon,1-\delta]$, and $I_\plus$ is continuous there. So the Riemann
sum gives $\sum_kG_N(k)\sim NT_\plus(N)\int_{1-\varepsilon}^{1-\delta}I_\plus(v)\,dv$. The points
$(1-\delta)N<k<N$ have $L_k\ge N/2$ for small $\delta$, so $G_N(k)\le C\widehat G_N(k)\le CT_\plus(N)$
by~\eqref{eq:G-size}, and by the additive form again they contribute at most
$C\delta NT_\plus(N)$. The atom at $k=N$ is $T_\plus(N)=o(NT_\plus(N))$. Let $\delta\to0$. With
$\ell=(v-p)/(1-p)$, $dv=(1-p)\,d\ell$ and $(1-p)/\mu_\minus=1/s$,
\[
 \int_{1-\varepsilon}^{1}I_\plus(v)\,dv=\frac1s\int_{\ell_\varepsilon}^{1}\bigl[\alpha(1-\ell)\ell^{-\alpha-1}+\ell^{-\alpha}\bigr]d\ell
 =\frac1s\Bigl[-(1-\ell)\ell^{-\alpha}\Bigr]_{\ell_\varepsilon}^{1}=\frac{(1-\ell_\varepsilon)\ell_\varepsilon^{-\alpha}}{s}
 =\frac{\varepsilon}{\mu_\minus}\,\ell_\varepsilon^{-\alpha},
\]
since $(1-\ell)\ell^{-\alpha}$ has derivative $-\alpha(1-\ell)\ell^{-\alpha-1}-\ell^{-\alpha}$ and
$1-\ell_\varepsilon=\varepsilon/(1-p)$. For the random-time start, Corollary~\ref{cor:stationary}(b)
gives the ratio form $P_N(k)\sim G^{\mathrm{st}}_N(k)$ uniformly on the zone with $k<N$, with
$G^{\mathrm{st}}_N(\lfloor vN\rfloor)/T_\plus(N)\to I_\plus(v)+\ell^{-\alpha}/\mu_\minus$ uniformly on
compacts, and $G^{\mathrm{st}}_N(k)\le CT_\plus(N)$ on $L_k\ge N/2$. The same argument adds
$\int_{1-\varepsilon}^1\ell^{-\alpha}\mu_\minus^{-1}\,dv=(\ell_\varepsilon^{1-\alpha}-1)/(s(\alpha-1))$ to
the integral and the atom $Q_\plus(N)\sim NT_\plus(N)/(s(\alpha-1))$ to the sum, which together
give the second term.
\end{proof}

Part (b) is the design rule. Starting at a regeneration, that is on an empty system, makes
the no-switch probability smaller by a factor of order $N$ than starting at an arbitrary
time, and by Corollary~\ref{cor:stationary}(b) it also removes the complete initial cycle
from the big-jump kernel, which is the second term in (a). The plus and minus statements
differ only by the two loads, so the heavier arm, the one with the larger load, is the one
whose cycle is more likely to swallow the experiment.

\begin{corollary}[the most likely exposure]\label{cor:exposure-mode}
\begin{enumerate}[(a)]
\item For a fresh start of either arm, the most likely value of $A_N$ is interior for all
large $N$, within $O(B_N)$ of $pN$.
\item For a random-time start, and for an aged start with $t_a/N\to\theta\in(0,\infty]$, the
single most likely value of $A_N$ is an endpoint, that is no switch, for all large $N$ when
$\alpha<\varphi=(1+\sqrt5)/2$, and is interior for all large $N$ when $\alpha>\varphi$. The
same holds for an experiment started at a typical time of a system run in control only,
with the cycle in progress kept in control.
\item If $Y$ also satisfies (LD2), the start is fresh or fair, the no-tie condition of
Theorem~\ref{thm:ends} holds at both ends and each end profile $u_\mp(j)$, $j\ge1$, has a single
local maximum, then $A_N$ has exactly 3 modes for all large $N$, the bulk mode and the two
no-switch outcomes. The last condition holds when $f_\pm$ is strictly log-convex.
\end{enumerate}
\end{corollary}

\begin{proof}
(a) and the random-time and aged cases of (b) are Theorem~\ref{thm:mode}(d) and
Corollary~\ref{cor:aged-mode}, with $\alpha_\plus=\alpha_\minus=\alpha$. For the control-only
start the initial run is $\nu_\minus=T_\minus/\mu_\minus$ and $\nu_\plus=0$, so
$B_n\Prob(D=n)=B_nT_\minus(n)/\mu_\minus\to0$ and $Q_\minus(N)=\mu_\minus^{-1}\sum_{d\ge N}T_\minus(d)\in\RV_{1-\alpha}$,
and Theorem~\ref{thm:mode}(c) gives the same threshold. (c) By
Proposition~\ref{prop:cycle}, the cycle laws satisfy (SPAN), (LR), (LD2) and (EQ), so
Theorem~\ref{thm:mode-count} and Corollary~\ref{cor:three-modes} apply, and
Corollary~\ref{cor:logconvex-three} covers the log-convex case.
\end{proof}

Four points of (b) deserve care. The threshold is not an equivalence at $\alpha=\varphi$,
where slowly varying factors decide. For the cycle laws of the queue the tails are pure
powers up to the factor of $T_Y$, and the ratio $B_NQ_\sgn(N)/H_N$ tends to a constant that
we have not computed. The theorem compares each endpoint atom with the largest interior
value, not the sum of the two atoms with it. The statement holds for all large $N$, not for
every $N$. With different service laws in the two arms the indices may differ, and the
criterion is then $\alpha_\sgn-1<1/\alpha_*$ for each arm separately, by
Theorem~\ref{thm:mode}(c). For (c), strict log-convexity of the total-progeny law is
checked only numerically, for the zeta service law, and we do not know a proof.
Section~\ref{sec:verify} reports the checks of this section against the exact law, computed
from~\eqref{eq:cyclic} and the recursion~\eqref{eq:run-recursion}, and a simulation of the
queue itself.

\ifarxivbuild
\begin{remark}[switching after several cycles]\label{rem:geometric-cycles}
A design that switches arm after a geometric number $G$ of cycles, with
$\Prob(G=n)=g(1-g)^{n-1}$, has runs $R=\sum_{i\le G}C_i$. By Asmussen, Foss and
Korshunov~\cite[Theorem~2(i) and Remark~3]{afk2003}, $f_R(n)\sim\E[G]f_C(n)$, so (LR) and (SPAN)
carry over to $R$ and Corollaries~\ref{cor:exposure-bulk} and~\ref{cor:exposure-tail} hold with
the run laws $f_{R_\pm}$. We have not checked (LD) and (LD2) for $R$.
\end{remark}
\fi
